
\begin{frame}{Anhang: Kompetenzen als theoretischer Kontext}
\centering
\resizebox{.74\textwidth}{!}{% Adapted from the user's updated framework. The key is on the next slide.
\begingroup
\tikzset{
original/.style={line width=1.1pt},
relation/.style={line width=1.5pt,{Triangle[length=2mm]}-{Triangle[length=2mm]}},
labelbox/.style={draw=black,line width=.4pt,fill=white,inner sep=3pt,align=left,font=\fontsize{10}{11}\selectfont},
maintext/.style={align=center,font=\bfseries\fontsize{11}{12}\selectfont,inner sep=0pt},
badge/.style={circle,draw=competencygreen,fill=white,text=competencygreen,minimum size=4.5mm,inner sep=0pt,font=\bfseries\small},
feedback/.style={draw=competencygreen,line width=1.1pt,-{Triangle[length=2mm]}}
}
\begin{tikzpicture}[x=.01cm,y=-.01cm]
% Learner, scientific model, and phenomenon.
\draw[original] (85,686)
  .. controls (60,660) and (46,626) .. (46,591)
  .. controls (46,504) and (122,433) .. (216,433)
  .. controls (310,433) and (385,504) .. (385,591)
  .. controls (385,626) and (374,659) .. (350,686)
  .. controls (373,714) and (385,749) .. (385,785)
  .. controls (385,872) and (310,942) .. (216,942)
  .. controls (122,942) and (46,872) .. (46,785)
  .. controls (46,748) and (60,713) .. (85,686);
\draw[original,densely dotted] (85,686)--(350,686);
\draw[original] (650,316)
  .. controls (637,295) and (630,272) .. (630,248)
  .. controls (630,161) and (706,90) .. (800,90)
  .. controls (894,90) and (970,161) .. (970,248)
  .. controls (970,273) and (964,296) .. (952,316)
  .. controls (963,338) and (969,360) .. (969,384)
  .. controls (969,471) and (893,541) .. (800,541)
  .. controls (706,541) and (629,471) .. (629,384)
  .. controls (629,359) and (637,335) .. (650,316);
\draw[original,densely dotted] (650,316)--(952,316);
\draw[original] (1180,432) rectangle (1518,942);
\draw[original] (1180,686)--(1518,686);

% Relations. Orange labels describe scientific processes, not assessed skills.
\draw[relation] (412,569)--(1141,569);
\draw[relation] (412,808)--(1141,808);
\draw[relation,alignmentblue] (385,571)--(632,246);
\draw[relation,alignmentblue] (385,786)--(637,466);
\draw[relation,modellingorange] (970,246)--(1180,578);
\draw[relation,modellingorange] (950,470)--(1180,812);

\node[maintext] at (216,380) {Lernende\\Mentale Modelle};
\node[maintext] at (797,69) {Physikalisches Modell};
\node[maintext] at (1348,407) {Reales Phänomen};
\node[maintext] at (217,561) {Gestalt};
\node[maintext] at (217,815) {Funktionalität};
\node[maintext] at (800,204) {Ausgewählte\\Formen und\\Strukturen};
\node[maintext] at (800,402) {Ausgewählte\\Funktionen};
\node[maintext] at (1348,561) {Sinnliche\\Muster};
\node[maintext] at (1348,815) {Funktionen};
\node[labelbox,text=alignmentblue] at (492,430) {Strukturelles\\Alignment};
\node[labelbox,text=alignmentblue] at (504,694) {Funktionales\\Alignment};
\node[labelbox,text=modellingorange] at (1075,385) {Strukturelle\\Modellierung};
\node[labelbox,text=modellingorange] at (1075,663) {Funktionale\\Modellierung};
\node[labelbox,align=center] at (777,605) {Fidelity of Gestalt};
\node[labelbox,align=center] at (792,855) {Functional Fidelity};

% Green markers denote the learner's activities at each relation.
% 3 and 4 mean mapping model to phenomenon, not constructing scientific models.
\node[badge] at (606,410) {1};
\node[badge] at (620,675) {2};
\node[badge] at (1138,319) {3};
\node[badge] at (1138,597) {4};
\node[badge] at (946,605) {6};
\node[badge] at (960,855) {6};

% Higher-order competency 5: compact upper-left annotation and smooth return.
\draw[feedback,line cap=round,line join=round] (650,165)
  .. controls (540,-20) and (20,-20) .. (20,185)
  .. controls (20,290) and (216,285) .. (216,340);
\node[anchor=north west,inner sep=3pt,align=left,text width=5cm,
  text=competencygreen,font=\fontsize{7.5}{9.5}\selectfont] at (75,80)
  {\textbf{Modellauswahl und -transformation}\\[3pt]
   Modelle und Repräsentationen zielgerichtet\\
   auswählen, koordinieren und für das\\
   eigene Verständnis nutzen};
\node[badge] at (294,37) {5};
% Evidence-based revision returns to both components of the mental model.
\draw[feedback] (1350,950)
  .. controls (1350,1006) and (1310,1020) .. (1230,1020)
  -- (80,1020)
  .. controls (-62,1020) and (-65,956) .. (-65,861)
  -- (-65,648)
  .. controls (-65,590) and (-28,590) .. (44,590);
\draw[feedback] (-65,850)
  .. controls (-65,793) and (-20,793) .. (43,793);
\node[fill=white,inner sep=4pt,align=center,text=competencygreen,
  font=\fontsize{8.2}{10.6}\selectfont] at (740,1020)
  {\textbf{Kompetenz der Lernenden: Aktualisierung des mentalen Modells}\\
   Testen und bewerten; Gestalt und Funktionalität anhand von Evidenz überarbeiten};
\node[badge] at (210,1020) {7};
\end{tikzpicture}
\endgroup
}
% \source{Grüne Markierungen: Aktivitäten der Lernenden. Orange Relationen: wissenschaftliche Modellierungsprozesse.}
\note{Das Rahmenmodell unterscheidet das mentale Modell der Lernenden, das physikalische Modell und das reale Phänomen. Fidelity beschreibt die Überzeugungen der Lernenden; Alignment beschreibt die Übereinstimmung mit dem Expertenkonsens. Orangefarbene Relationen bezeichnen wissenschaftliche Modellierungsprozesse bei der Entwicklung von Modellen und wissenschaftlichem Konsens. Die grünen Markierungen 3 und 4 bezeichnen das Erkennen und Vergleichen von Korrespondenzen entlang dieser Relationen durch die Lernenden, nicht das Konstruieren wissenschaftlicher Modelle. Rückpfeil 5 verbindet Modellauswahl und -transformation mit dem Verständnis der Lernenden. Rückpfeil 7 verbindet Testen und Bewerten mit der evidenzbasierten Überarbeitung sowohl der Gestalt als auch der Funktionalität. Die sieben Kompetenzen bilden den theoretischen Kontext und sind keine gemessenen Zielgrößen dieser Studie.}
\end{frame}

\begin{frame}{Anhang: Beispiele für Modellierungspraktiken}
\footnotesize\renewcommand{\arraystretch}{1.04}
\begin{tabularx}{\textwidth}{@{}P{.025\textwidth}P{.32\textwidth}Y@{}}
\toprule
& \textbf{Kompetenz der Lernenden} & \textbf{Illustratives Beispiel (keine zusätzliche Aufgabe)}\\
\midrule
1 & Strukturelle Zuordnung:\newline mental--physikalisch & Die vorgestellte Polarisationsellipse mit der ausgewählten Struktur des physikalischen Modells vergleichen.\\
2 & Funktionale Zuordnung:\newline mental--physikalisch & Die eigene Erklärung von Phaseneffekten mit den funktionalen Zusammenhängen des physikalischen Modells vergleichen.\\
3 & Strukturelle Modellierung:\newline Zuordnung Physik--Welt & Die dargestellte Orientierung eines Polarisators der Transmissionsachse eines realen Polarisators zuordnen.\\
4 & Funktionale Modellierung:\newline Zuordnung Physik--Welt & Die vom Modell vorhergesagte Transmission mit Änderungen der realen Lichtintensität in Beziehung setzen.\\
5 & Modellauswahl und -transformation & Die Jones-Notation zur Berechnung einer Transformation und die Kugel zu deren geometrischer Interpretation nutzen.\\
6 & Metakognitive Reflexion & Überzeugungen über den visuellen Realismus und die funktionale Angemessenheit der Kugel einschließlich ihrer Abstraktionen und Grenzen reflektieren.\\
7 & Evidenzbasierte Überarbeitung & Eine Diskrepanz zwischen Vorhersage und Applet-Feedback nutzen, um Gestalt, Funktionalität oder beides zu revidieren.\\
\bottomrule
\end{tabularx}
% \source{Beispiele adaptiert an das aktualisierte Rahmenmodell. Modellkompetenzen werden in dieser Studie nicht erhoben.}
\note{Die Zeilen 1 und 2 betreffen Korrespondenzen zwischen dem mentalen Modell der Lernenden und dem physikalischen Modell. Die Zeilen 3 und 4 betreffen die Zuordnung zwischen physikalischem Modell und Phänomen durch die Lernenden. Ihre Bezeichnungen folgen dem vorgegebenen Rahmenmodell; die Beispiele implizieren nicht, dass Lernende einen wissenschaftlichen Konsens konstruieren oder dass Applet-Feedback ein unabhängiges physikalisches Experiment darstellt. Zeile 5 umfasst die zielgerichtete Auswahl, Transformation und Koordination von Repräsentationen zur Förderung des eigenen Verständnisses. Zeile 6 betrifft die Reflexion über Fidelity-Überzeugungen, während das Experten-Alignment einen separaten Vergleich mit von Expert:innen ausgewählten Inhalten darstellt. Zeile 7 aktualisiert beide möglichen Komponenten des mentalen Modells anhand von Evidenz. Diese Beispiele erläutern das Rahmenmodell; es handelt sich nicht um zusätzliche Aktivitäten oder Interviewimpulse. Laut-Denken-Beobachtungen bleiben auf spontanes Schlussfolgern während der bestehenden Applet-Aufgaben beschränkt.}
\end{frame}
